% Motor: XeLaTeX. Sesión 31 - material docente UTP.
% Estructura docente de la Sesión 30.
\documentclass[aspectratio=169,10pt,c]{beamer}
% Reutiliza sin cambios la identidad visual UTP validada en la Sesion 08.
% Reutiliza sin cambios la identidad visual UTP validada en la Sesion 05.
% La segunda ruta permite compartir los assets comunes sin duplicar binarios.
% !TEX root = slides.tex
% Preambulo comun: estilo UTP, colores, fuentes, plantillas de beamer y
% comandos de apoyo didactico. Se carga desde slides.tex tras \documentclass.

% La opción "c" de beamer reparte el espacio libre 1:1.5 (arriba:abajo), por lo
% que el contenido queda por encima del centro. Igualamos ambos rellenos.
\makeatletter
\define@key{beamerframe}{c}[true]{%
  \beamer@frametopskip=0pt plus 1fill\relax%
  \beamer@framebottomskip=0pt plus 1fill\relax%
  \beamer@frametopskipautobreak=0pt plus .5\paperheight\relax%
  \beamer@framebottomskipautobreak=0pt plus .5\paperheight\relax%
}
\makeatother

\usepackage{graphicx}
\usepackage{tikz}
\usepackage{xcolor}
\usepackage{iftex}

\ifPDFTeX
  \usepackage[utf8]{inputenc}
  \usepackage[T1]{fontenc}
  \usepackage{lmodern}          % tipografia escalable para el respaldo con pdflatex
  \renewcommand{\familydefault}{\sfdefault}
\else
  \usepackage{fontspec}
  \IfFontExistsTF{Arial}
    {\setsansfont{Arial}[BoldFont={Arial Bold}]}
    {\setsansfont{Helvetica}[BoldFont={Helvetica Bold}]}
\fi

\usepackage[spanish,es-nodecimaldot]{babel}
\usepackage{microtype}
\usepackage{booktabs}
\usepackage{array}
\usetikzlibrary{arrows.meta,positioning,fit,calc,shapes.geometric}

\graphicspath{{assets/}}

% ---------- Identidad visual UTP ----------
\definecolor{utpRojo}{HTML}{C42B35}
\definecolor{utpNegro}{HTML}{111111}
\definecolor{utpGris}{HTML}{5F5F5F}
\definecolor{utpGrisMedio}{HTML}{9A9A9A}
\definecolor{utpGrisClaro}{HTML}{F1F1F1}

\usetheme{default}
\setbeamercolor{normal text}{fg=utpNegro,bg=white}
\setbeamercolor{structure}{fg=utpRojo}
\setbeamercolor{frametitle}{fg=utpNegro,bg=white}
\setbeamercolor{title}{fg=utpNegro}
\setbeamercolor{subtitle}{fg=utpNegro}
\setbeamercolor{author}{fg=utpNegro}
\setbeamercolor{date}{fg=utpNegro}
\setbeamercolor{block title}{fg=white,bg=utpNegro}
\setbeamercolor{block body}{fg=utpNegro,bg=utpGrisClaro}
\setbeamercolor{alertblock title}{fg=white,bg=utpRojo}
\setbeamercolor{alertblock body}{fg=utpNegro,bg=utpRojo!7}
\setbeamercolor{exampleblock title}{fg=white,bg=utpGris}
\setbeamercolor{exampleblock body}{fg=utpNegro,bg=utpGris!8}
\setbeamercolor{terminal}{fg=white,bg=utpNegro}

\setbeamerfont{title}{series=\bfseries,size=\fontsize{19.5}{22.5}\selectfont}
\setbeamerfont{subtitle}{size=\fontsize{11.5}{13.5}\selectfont}
\setbeamerfont{date}{series=\bfseries,size=\fontsize{11.5}{13.5}\selectfont}
\setbeamerfont{author}{size=\fontsize{10.5}{12.5}\selectfont}
\setbeamerfont{frametitle}{series=\bfseries,size=\large}
\setbeamerfont{block title}{series=\bfseries,size=\small}
\setbeamerfont{block body}{size=\normalsize}

\setbeamertemplate{navigation symbols}{}
\setbeamertemplate{headline}{}
\setbeamertemplate{blocks}[default]
\setbeamertemplate{itemize item}{\color{utpRojo}\small\raise.15ex\hbox{$\blacktriangleright$}}
\setbeamertemplate{itemize subitem}{\color{utpGris}\scriptsize$\bullet$}
\setbeamertemplate{enumerate item}{\color{utpRojo}\insertenumlabel.}
\setbeamertemplate{background canvas}{\color{white}\rule{\paperwidth}{\paperheight}}

% Título de diapositiva: negrita negra a la izquierda, logo UTP a la derecha
\setbeamertemplate{frametitle}{%
  \begin{tikzpicture}[remember picture,overlay]
    \node[anchor=north east,inner sep=0pt]
      at ([xshift=-7mm,yshift=-5.5mm]current page.north east)
      {\includegraphics[width=17mm]{utp-logo.png}};
  \end{tikzpicture}%
  \vspace{1.2mm}%
  \begin{minipage}[t]{.80\textwidth}\raggedright\insertframetitle\end{minipage}\par
  \vspace{1.1mm}{\color{utpRojo}\rule{\textwidth}{1.1pt}}
}

% Solo numeración discreta; la plantilla UTP no lleva pie de página
\setbeamertemplate{footline}{%
  \hfill{\color{utpGrisMedio}\scriptsize\insertframenumber}\hspace{6mm}\vspace{2.5mm}%
}

% Portada UTP
\setbeamertemplate{title page}{%
  \begin{tikzpicture}[remember picture,overlay]
    \node[anchor=center,inner sep=0pt] at (current page.center) {%
      \includegraphics[width=\paperwidth,height=\paperheight]{utp-cover-bg.jpeg}%
    };
    \node[anchor=north,align=center,text=utpNegro,inner sep=0pt,
      text width=145mm,font=\usebeamerfont{title}]
      at ([yshift=-25.2mm]current page.north) {\inserttitle};
    \node[anchor=north,align=center,text=utpNegro,inner sep=0pt,
      text width=142mm,font=\usebeamerfont{subtitle}]
      at ([yshift=-38.2mm]current page.north) {\insertsubtitle};
    \node[anchor=north,align=center,text=utpNegro,inner sep=0pt,
      font=\usebeamerfont{date}]
      at ([yshift=-50.4mm]current page.north) {\insertdate};
    \node[anchor=north,align=center,text=utpNegro,inner sep=0pt,
      font=\usebeamerfont{author}]
      at ([yshift=-58.4mm]current page.north) {\insertauthor};
  \end{tikzpicture}%
}

% ---------- Comandos de apoyo didáctico ----------
\newcommand{\concepto}[2]{%
  \begin{block}{#1}\large #2\end{block}}
\newcommand{\pregunta}[1]{%
  \begin{alertblock}{Decide y explica}#1\end{alertblock}}
\newcommand{\analogia}[1]{%
  \begin{exampleblock}{En palabras simples}\small #1\end{exampleblock}}
\newcommand{\enlarealidad}[1]{%
  \begin{block}{Cómo se ve en la realidad}\small #1\end{block}}
\newcommand{\ojo}[1]{%
  \begin{alertblock}{Ojo con esto}\small #1\end{alertblock}}
\newcommand{\apoyo}{{\normalfont\footnotesize\color{utpGris}~· material de apoyo}}
\newcommand{\marca}[1]{\tikz[baseline=(m.base)]{%
  \node[rounded corners=1.5pt,fill=utpGrisClaro,draw=utpGrisMedio,inner xsep=3pt,inner ysep=1.5pt,
        font=\footnotesize] (m) {#1};}}
\newenvironment{terminal}{%
  \begin{beamercolorbox}[wd=\linewidth,sep=5pt,rounded=true]{terminal}%
  \ttfamily\small\ignorespaces}{%
  \end{beamercolorbox}}

\tikzset{
  caja/.style={draw=utpNegro, very thick, rounded corners=1.5pt, fill=white,
    align=center, minimum height=0.72cm, inner sep=5pt, font=\small},
  cajaazul/.style={caja, draw=utpGris, fill=white},
  cajaverde/.style={caja, draw=utpNegro, fill=utpGrisClaro},
  cajanaranja/.style={caja, draw=utpRojo!75, fill=utpRojo!7},
  flecha/.style={-{Stealth[length=2.3mm]}, very thick, draw=utpNegro},
  flechavuelta/.style={{Stealth[length=2.3mm]}-{Stealth[length=2.3mm]}, very thick, draw=utpNegro}
}

\graphicspath{{../../Sesion05/slides/assets/}}



\graphicspath{{assets/}{../../Sesion05/slides/assets/}}
\usepackage{colortbl}
\usepackage{listings}
% Un archivo: las palabras clave en rojo.
\lstset{basicstyle=\ttfamily\scriptsize,backgroundcolor=\color{utpGrisClaro},
  frame=none,aboveskip=4pt,belowskip=4pt,showstringspaces=false,
  columns=fullflexible,keepspaces=true,breaklines=true,
  keywords={FROM,WORKDIR,COPY,RUN,CMD,runtime},
  keywordstyle=\color{utpRojo}\bfseries,
  commentstyle=\color{utpGris}\itshape,morecomment=[l]{\#},
  xleftmargin=5pt,xrightmargin=5pt,framesep=5pt}
% Una terminal: fondo oscuro, sin colorear palabras. Lo que va entre | | se resalta.
\lstdefinestyle{terminal}{basicstyle=\ttfamily\scriptsize\color{white},backgroundcolor=\color{utpNegro},
  keywords={},morecomment=[l]{\#},commentstyle=\color{utpGrisMedio},upquote=true,
  moredelim=[is][\color{yellow}\bfseries]{|}{|}}
\usepackage{upquote}
\usepackage{pifont}
\definecolor{okVerde}{HTML}{2E7D32}
% Un número en un círculo rojo, para unir lo que se necesita con su pieza.
\newcommand{\num}[1]{\textcolor{utpRojo}{\ding{#1}}}
\usetikzlibrary{decorations.pathreplacing,positioning,shadows}

\title{Integración de Servicios Cloud GCP}
\subtitle{App Engine, Cloud Run y pruebas de carga}
\date{Semana 16 - Sesión 31}
\author{Rel Guzman}

% Íconos simples para los diagramas.
\tikzset{
 caja/.style={draw=utpRojo,thick,rounded corners=4pt,align=center},
 zona/.style={draw=utpGrisMedio,dashed,rounded corners=4pt},
 flecha/.style={-latex,thick},
 pod/.style={draw=utpRojo,thick,rounded corners=3pt,fill=utpGrisClaro,align=center,font=\scriptsize\ttfamily,minimum width=1.2cm,minimum height=.5cm},
 version/.style={draw=utpRojo,thick,rounded corners=3pt,fill=utpRojo!30,align=center,font=\scriptsize\ttfamily,minimum width=1.9cm,minimum height=.6cm},
}

\begin{document}

{\setbeamertemplate{footline}{}\maketitle}

% 2
\begin{frame}{Revisión: un sistema que necesita Kubernetes}
\begin{center}
\begin{tikzpicture}[font=\small]
  \node[caja,dashed,text width=2.0cm,minimum height=1.0cm] (us) at (0,0) {\textbf{Clientes}\\\footnotesize lunes, 8 a.\,m.};
  \node[caja,fill=utpGrisClaro,text width=1.5cm,minimum height=.9cm] (web) at (4.3,0) {\textbf{web}};
  \node[caja,fill=utpRojo!30,text width=1.5cm,minimum height=.9cm] (api) at (6.6,0) {\textbf{api}};
  \node[zona,fit=(web)(api),inner sep=8pt,label={[font=\scriptsize,text=utpGris]below:una sola computadora}] (pc) {};
  \draw[flecha] (us) -- node[above,font=\scriptsize,align=center]{cientos\\de pedidos} (pc.west);
  \draw[flecha] (web) -- (api);
  \node[align=left,font=\footnotesize,anchor=west,text width=4.1cm] at (8.4,0) {\textcolor{utpRojo}{\ding{55}} Si la computadora se apaga, se cae todo.\\[3pt] \textcolor{utpRojo}{\ding{55}} Una sola copia de la API no alcanza.\\[3pt] \textcolor{utpRojo}{\ding{55}} Si la API se cae, alguien la levanta a mano.};
\end{tikzpicture}
\end{center}
\vspace{3mm}
\textbf{\small Lo que necesita Café Nube}\par\smallskip
\begin{columns}[T,onlytextwidth]
\column{.235\textwidth}
\begin{block}{\ding{202}}
\small Varias máquinas que trabajen como una
\end{block}
\column{.235\textwidth}
\begin{block}{\ding{203}}
\small Varias copias de la API
\end{block}
\column{.235\textwidth}
\begin{block}{\ding{204}}
\small Que las copias se reparen solas
\end{block}
\column{.235\textwidth}
\begin{block}{\ding{205}}
\small Un nombre fijo para encontrar las copias
\end{block}
\end{columns}
\end{frame}

% 3
\begin{frame}{Revisión: qué pieza de Kubernetes resuelve cada cosa}
\begin{center}
\begin{tikzpicture}[font=\small,yscale=.85]
  % lo que tú escribes y quien lo recibe
  \node[caja,dashed,text width=2.1cm,minimum height=.9cm] (man) at (0.6,2.1) {\textbf{Manifiesto}\\\scriptsize\texttt{deployment.yaml}};
  \node[caja,fill=utpGrisClaro,text width=3.6cm,minimum height=.8cm] (mae) at (5.7,2.1) {\textbf{Nodo maestro}\\\scriptsize decide dónde va cada pod};
  \draw[flecha] (man) -- node[above,font=\scriptsize\ttfamily,align=center]{kubectl\\apply} (mae);
  % las máquinas, con un pod de la API en cada una
  \foreach \i/\x in {1/3.6,2/5.7,3/7.8} {
    \node[caja,text width=1.6cm,minimum height=1.1cm] (n\i) at (\x,0) {};
    \node[font=\scriptsize\bfseries,anchor=north] at (n\i.north) {nodo \i};
    \node[pod,fill=utpRojo!30] (p\i) at (\x,-.18) {api};
  }
  % las copias que cuida el Deployment
  \draw[decorate,decoration={brace,amplitude=5pt},thick] (2.75,.75) -- node[above=5pt,font=\scriptsize,fill=white,inner sep=1pt]{\textbf{Deployment} \texttt{api}: 3 réplicas} (8.65,.75);
  \draw[flecha,utpGrisMedio] (mae.south) -- (5.7,1.25);
  % el nombre fijo que reparte los pedidos
  \node[caja,text width=1.8cm,minimum height=.9cm] (svc) at (0.6,0) {\textbf{Service}\\\scriptsize\texttt{api}};
  \draw[flecha] (svc) -- (p1);
  \draw[flecha] (svc.south east) to[out=-10,in=-90,looseness=.5] (p2.south);
  \draw[flecha] (svc.south) to[out=-10,in=-90,looseness=.5] (p3.south);
  \node[zona,fit=(mae)(n1)(n3),inner ysep=9pt,inner xsep=6pt,label={[font=\scriptsize,text=utpGris]above:clúster}] {};
\end{tikzpicture}
\end{center}
\vspace{-7mm}
\renewcommand{\arraystretch}{1.2}
\small
\begin{tabular}{@{}>{\raggedright\arraybackslash}p{4.35cm}>{\raggedright\arraybackslash}p{4.35cm}>{\raggedright\arraybackslash}p{4.35cm}@{}}
\num{202} \textbf{Clúster}: varias máquinas que trabajan como una & \textbf{Nodo}: cada máquina. El \textbf{maestro} reparte los pods & \textbf{Pod}: una copia de tu app corriendo, con su contenedor y su IP. Vive dentro de un nodo; un nodo puede tener varios\\
\num{203}\,\num{204} \textbf{Deployment}: mantiene las copias; si falta una, crea otra & \num{205} \textbf{Service}: el nombre fijo que reparte los pedidos & \textbf{Manifiesto}: un YAML que describe lo que quieres\\
\end{tabular}
\end{frame}

% 4
\begin{frame}{Revisión: Café Nube en Kubernetes}
\begin{center}
\begin{tikzpicture}[font=\small,
    paso/.style={caja,fill=white,text width=1.45cm,minimum height=1.1cm},
    nota/.style={font=\scriptsize,text=utpGris,anchor=north,text width=2.2cm,align=center}]
  % un pedido de izquierda a derecha; las cartas apiladas son las copias
  \node[paso,dashed] (us) at (0,0) {\textbf{Cliente}\\\scriptsize\texttt{/api/menu}};
  \node[paso] (sw) at (2.85,0) {\textbf{Service}\\\scriptsize\texttt{web}};
  \node[paso,fill=utpGrisClaro,copy shadow] (web) at (5.7,0) {\textbf{pod web}\\\scriptsize ×2};
  \node[paso] (sa) at (8.55,0) {\textbf{Service}\\\scriptsize\texttt{api:8080}};
  \node[paso,fill=utpRojo!30,double copy shadow] (api) at (11.4,0) {\textbf{pod api}\\\scriptsize ×3};
  \draw[flecha] (us) -- (sw);
  \draw[flecha] (sw) -- node[above,font=\scriptsize,align=center]{a una\\de las 2} (web);
  \draw[flecha] (web) -- (sa);
  \draw[flecha] (sa) -- node[above,font=\scriptsize,align=center]{a una\\de las 3} (api);
  \node[nota] (c1) at (2.85,-.75) {nombre fijo\\de la página};
  \node[nota] (c2) at (5.7,-.75) {2 copias de la página\\(Deployment \texttt{web})};
  \node[nota] (c3) at (8.55,-.75) {nombre fijo\\de la API};
  \node[nota] (c4) at (11.4,-.75) {3 copias de la API\\(Deployment \texttt{api})};
  \node[zona,fit=(sw)(api)(c1)(c4),inner sep=6pt,label={[font=\scriptsize,text=utpGris]above:clúster}] {};
\end{tikzpicture}
\end{center}
{\small \textbf{Los nodos:} el \textbf{nodo maestro} decide dónde va cada pod. \textbf{Nodo 1}: \texttt{api} y \texttt{web}. \textbf{Nodo 2}: \texttt{api} y \texttt{web}. \textbf{Nodo 3}: \texttt{api}.}
\vspace{2mm}
\begin{columns}[T,onlytextwidth]
\column{.48\textwidth}
\begin{block}{Si se borra un pod \texttt{api}}
\small El Deployment ve 2 de 3 y \textbf{crea otro}, con otra IP. El Service lo suma solo: \texttt{web} sigue pidiendo a \texttt{api}.
\end{block}
\column{.48\textwidth}
\begin{block}{Si se apaga un nodo}
\small Kubernetes vuelve a arrancar sus pods \textbf{en otro nodo}. Café Nube sigue respondiendo.
\end{block}
\end{columns}
\end{frame}

% 5
\begin{frame}{Revisión: Kubernetes en Google Cloud, GKE}
\begin{center}
\begin{tikzpicture}[font=\small]
  \node[caja,fill=utpGrisClaro,text width=2.3cm,minimum height=1.0cm] (cb) at (0,0) {\textbf{Cloud Build}\\\footnotesize construye la imagen};
  \node[caja,fill=utpRojo!30,text width=2.3cm,minimum height=1.0cm] (ar) at (3.8,0) {\textbf{Artifact Registry}\\\footnotesize guarda la imagen};
  \node[caja,draw=okVerde,fill=okVerde!12,text width=2.3cm,minimum height=1.0cm] (gke) at (7.6,0) {\textbf{GKE}\\\footnotesize la corre en pods};
  \draw[flecha] (cb) -- (ar);
  \draw[flecha] (ar) -- (gke);
\end{tikzpicture}
\end{center}
\vspace{2mm}
\begin{center}
\begin{minipage}{.72\textwidth}
\begin{block}{GKE}
\small Kubernetes administrado por Google. Google maneja el nodo maestro; los nodos son VMs de Compute Engine. Tiene los mismos pods, Deployments y Services, y usa los \textbf{mismos manifiestos}.
\end{block}
\end{minipage}
\end{center}
\end{frame}

% 6
\begin{frame}[c]
\begin{block}{Logro de la sesión}
Al finalizar, podrás implementar la API de Flask de Café Nube en \textbf{App Engine} y en \textbf{Cloud Run}, publicar una \textbf{versión nueva sin tráfico} y repartir el tráfico entre versiones, y medir con \textbf{Apache Bench} cuántos pedidos por segundo responde y si falla alguno.
\end{block}
\end{frame}

% 7
\begin{frame}{Saberes previos}
\begin{columns}[T,onlytextwidth]
\column{.46\textwidth}
\begin{block}{¿Qué es…?}
\begin{itemize}\setlength\itemsep{8pt}
  \item ¿Qué es \textbf{Flask}?
  \item ¿Qué es una \textbf{imagen} de contenedor?
  \item ¿Qué es un \textbf{Dockerfile}?
  \item ¿Qué es un \textbf{puerto}?
  \item ¿Qué es una \textbf{variable de entorno}?
\end{itemize}
\end{block}
\column{.50\textwidth}
\begin{block}{¿Qué hace…?}
\begin{itemize}\setlength\itemsep{8pt}
  \item En la API de Café Nube, ¿qué hace \texttt{@app.get("/menu")}?
  \item ¿Qué archivo dice qué bibliotecas instala \texttt{pip}, como \texttt{flask}?
  \item Si la API corre en una VM y nadie la visita en toda la noche, ¿la VM sigue encendida? ¿Sigue cobrando?
\end{itemize}
\end{block}
\end{columns}
\end{frame}

% 8
\begin{frame}[c]{Utilidad de esta sesión}
\begin{center}
\begin{minipage}{.84\textwidth}
{\large Con Kubernetes administras un \textbf{clúster}. Muchas apps no lo necesitan: en \textbf{App Engine} y \textbf{Cloud Run} subes tu código o tu imagen, y Google pone las máquinas y te da una URL. Antes de abrirla al público, una \textbf{prueba de carga} te dice cuántos pedidos aguanta.}
\end{minipage}
\end{center}
\end{frame}

% 9
\begin{frame}{Sin servidores: ¿qué haces tú y qué hace Google?}
\renewcommand{\arraystretch}{1.25}
\small
\begin{center}
\begin{tabular}{@{}>{\raggedright\arraybackslash}p{4.3cm}>{\raggedright\arraybackslash}p{3.9cm}>{\raggedright\arraybackslash}p{4.4cm}@{}}
\toprule
 & \textbf{En una VM} & \textbf{Sin servidores}\\
\midrule
Crear y actualizar la máquina & Tú & Google\\
Instalar Python y Flask & Tú, con \texttt{pip} & Google, con tu \texttt{requirements.txt}\\
Más copias si llegan más pedidos & Tú, a mano & \textcolor{okVerde}{\textbf{Google las crea y las quita solo}}\\
Toda la noche sin visitas & Sigue encendida y cobra & \textcolor{okVerde}{\textbf{Apaga todas las copias y no cobra por ellas}}\\
Una URL con HTTPS & Tú & Google te la da\\
\bottomrule
\end{tabular}
\end{center}
\vspace{2mm}
\begin{block}{Sin servidores (\textit{serverless})}
\small Sí hay servidores, pero \textbf{no los administras tú}. Subes tu código o tu imagen; Google crea las copias, llamadas \textbf{instancias}, según los pedidos que lleguen. Hoy vemos dos: \textbf{App Engine} y \textbf{Cloud Run}.
\end{block}
\end{frame}

% 10
\begin{frame}[fragile]{App Engine: subes tu código}
\begin{columns}[T,onlytextwidth]
\column{.40\textwidth}
\textbf{\scriptsize\texttt{main.py}}
\begin{lstlisting}
from flask import Flask, jsonify
app = Flask(__name__)
MENU = [...]

@app.get("/menu")
def menu():
    return jsonify(MENU)
\end{lstlisting}
\textbf{\scriptsize\texttt{requirements.txt}}
\begin{lstlisting}
flask==3.1.0
\end{lstlisting}
\textbf{\scriptsize\texttt{app.yaml}}
\begin{lstlisting}
runtime: python312
\end{lstlisting}
\column{.56\textwidth}
\begin{block}{Tres archivos}
\small \texttt{main.py}: tu API.\\ \texttt{requirements.txt}: las bibliotecas.\\ \texttt{app.yaml}: qué Python usar.
\end{block}
\begin{lstlisting}[style=terminal]
gcloud app create --region=us-central
gcloud app deploy
...
Deployed service [default] to
  [https://PROYECTO.uc.r.appspot.com]
\end{lstlisting}
\begin{exampleblock}{Y App Engine hace el resto}
\small Instala las bibliotecas, arranca tu API y te da una URL con HTTPS.
\end{exampleblock}
\end{columns}
\end{frame}

% 11
\begin{frame}[fragile]{Cloud Run: subes un contenedor}
\begin{center}
\begin{tikzpicture}[font=\small]
  \node[caja,dashed,text width=2.0cm,minimum height=1.1cm] (cod) at (0,0) {\textbf{Tu código}\\\footnotesize\texttt{main.py}, Dockerfile};
  \node[caja,fill=utpGrisClaro,text width=2.3cm,minimum height=1.1cm] (cb) at (3.6,0) {\textbf{Cloud Build}\\\footnotesize construye la imagen};
  \node[caja,fill=utpRojo!30,text width=2.3cm,minimum height=1.1cm] (ar) at (7.4,0) {\textbf{Artifact Registry}\\\footnotesize guarda la imagen};
  \node[caja,draw=okVerde,fill=okVerde!12,text width=2.3cm,minimum height=1.1cm] (run) at (11.2,0) {\textbf{Cloud Run}\\\footnotesize la corre y da una URL};
  \draw[flecha] (cod) -- (cb);
  \draw[flecha] (cb) -- (ar);
  \draw[flecha] (ar) -- (run);
\end{tikzpicture}
\end{center}
\vspace{1mm}
\begin{columns}[T,onlytextwidth]
\column{.53\textwidth}
\textbf{\scriptsize\texttt{Dockerfile}}
\begin{lstlisting}
FROM python:3.12-slim
WORKDIR /app
COPY requirements.txt .
RUN pip install -r requirements.txt
COPY . .
CMD exec gunicorn --bind :$PORT --threads 8 main:app
\end{lstlisting}
{\scriptsize \texttt{requirements.txt} suma \texttt{gunicorn}: aquí lo instalas tú. \mbox{\texttt{--threads 8}}: atiende hasta 8 pedidos a la vez.}
\column{.43\textwidth}
\begin{block}{El puerto}
\small Cloud Run le dice al contenedor en qué puerto escuchar con la variable \texttt{PORT}: \textbf{8080}. Por eso \texttt{--bind :\$PORT}.
\end{block}
\begin{block}{Igual que en la Sesión 30}
\small Cloud Build y Artifact Registry. Lo nuevo: \textbf{no hay clúster}. Le das la imagen y te devuelve una URL que termina en \texttt{run.app}.
\end{block}
\end{columns}
\end{frame}

% 12
\begin{frame}[fragile]{Cloud Run: crear el servicio}
\begin{lstlisting}[style=terminal,basicstyle=\ttfamily\footnotesize\color{white}]
IMAGEN=us-central1-docker.pkg.dev/PROYECTO/cafe/cafe-api
gcloud builds submit --tag $IMAGEN:v1 .
gcloud run deploy cafe-api --image $IMAGEN:v1 --region us-central1 \
    --allow-unauthenticated --max-instances 3
...
Service URL: https://cafe-api-123456789.us-central1.run.app
\end{lstlisting}
\vspace{1mm}
\renewcommand{\arraystretch}{1.2}
\small
\begin{tabular}{@{}>{\raggedright\arraybackslash}p{4.1cm}>{\raggedright\arraybackslash}p{8.7cm}@{}}
\toprule
\textbf{Opción} & \textbf{Qué decide}\\
\midrule
\texttt{--region} & Dónde corre. La imagen va en la misma región.\\
\texttt{--allow-unauthenticated} & \textbf{Acceso público}: cualquiera con la URL puede llamar a la API. Si pides autenticación, el navegador recibe \textbf{403 Forbidden}.\\
\texttt{--max-instances 3} & Cuántas copias puede crear como máximo: limita cuánto crece y cuánto cuesta.\\
\texttt{--set-env-vars} & Variables de entorno, como \texttt{-e} en Docker.\\
\bottomrule
\end{tabular}
\vspace{1mm}

{\scriptsize Lo mismo se hace en la consola: \textbf{Cloud Run > Implementar contenedor}.}
\end{frame}

% 13
\begin{frame}{Cloud Run: revisiones y tráfico}
\begin{center}
\begin{tikzpicture}[font=\small]
  \node[caja,fill=utpGrisClaro,text width=2.4cm,minimum height=1.0cm] (svc) at (0,0) {\textbf{Servicio}\\\footnotesize\texttt{cafe-api}};
  \node[version,minimum width=3.2cm] (r1) at (6,.7) {cafe-api-00001 (v1)};
  \node[version,minimum width=3.2cm,fill=okVerde!12,draw=okVerde] (r2) at (6,-.7) {cafe-api-00002 (v2)};
  \draw[flecha] (svc) -- node[above,sloped,font=\scriptsize]{90\%} (r1);
  \draw[flecha] (svc) -- node[below,sloped,font=\scriptsize]{10\%} (r2);
  \node[font=\scriptsize,text=utpGris,anchor=west,text width=3.4cm] at (8.0,-.7) {la nueva recibe un poco: si falla, pocos lo notan};
\end{tikzpicture}
\end{center}
\vspace{1mm}
\begin{columns}[T,onlytextwidth]
\column{.32\textwidth}
\begin{block}{Revisión}
\small Cada cambio de imagen o de configuración crea una \textbf{revisión} nueva. Las anteriores no cambian.
\end{block}
\column{.32\textwidth}
\begin{block}{Sin tráfico}
\small La revisión nueva arranca y se comprueba que responde, pero recibe \textbf{0\%}: \texttt{--no-traffic}.
\end{block}
\column{.32\textwidth}
\begin{block}{Canary}
\small Le das un poco de tráfico, como 10\%, para ver si falla con usuarios reales. Si va bien, 100\%.
\end{block}
\end{columns}
\end{frame}

% 14
\begin{frame}{App Engine y Cloud Run, lado a lado}
\renewcommand{\arraystretch}{1.3}
\small
\begin{center}
\begin{tabular}{@{}>{\raggedright\arraybackslash}p{3.9cm}>{\raggedright\arraybackslash}p{4.3cm}>{\raggedright\arraybackslash}p{4.3cm}@{}}
\toprule
 & \textbf{App Engine} & \textbf{Cloud Run}\\
\midrule
¿Qué subes? & Tu código & Una imagen de contenedor\\
¿Qué archivo dice cómo correrlo? & \texttt{app.yaml} & \texttt{Dockerfile}\\
¿Cuántas por proyecto? & Una aplicación & Muchos servicios\\
La URL termina en & \texttt{appspot.com} & \texttt{run.app}\\
Sin visitas & \multicolumn{2}{c}{\textcolor{okVerde}{\textbf{Las dos bajan a cero instancias}}}\\
\bottomrule
\end{tabular}
\end{center}
\vspace{2mm}
\begin{block}{¿Cuál uso?}
\small Si tu app es Python y no quieres escribir un Dockerfile, \textbf{App Engine}. Si ya tienes la imagen, como Café Nube desde la Sesión 28, \textbf{Cloud Run}.
\end{block}
\end{frame}

% 15
\begin{frame}[c]
\begin{center}
{\Huge\bfseries\color{utpRojo} Pruebas de carga}
\end{center}
\end{frame}

% 16
\begin{frame}{Prueba de carga: ¿cuánto aguanta mi API?}
\begin{columns}[T,onlytextwidth]
\column{.44\textwidth}
\begin{block}{Prueba de carga}
\small Mandar \textbf{muchos pedidos a la vez} a tu servicio, a propósito, para medir cómo responde \textbf{antes} de que lleguen los usuarios.
\end{block}
\begin{exampleblock}{Tres preguntas}
\small
1. ¿Cuántos pedidos por segundo responde?\\
2. ¿Cuánto espera cada usuario?\\
3. ¿Falla alguno?
\end{exampleblock}
\begin{alertblock}{Ojo}
\small Solo a \textbf{tus} servicios: a uno ajeno es un ataque. Y más pedidos son más instancias: más costo.
\end{alertblock}
\column{.52\textwidth}
\renewcommand{\arraystretch}{1.2}
\textbf{\small Algunas herramientas}\par\smallskip
\small
\begin{tabular}{@{}>{\raggedright\arraybackslash}p{2.2cm}>{\raggedright\arraybackslash}p{4.6cm}@{}}
\toprule
\textbf{Herramienta} & \textbf{Cómo se usa}\\
\midrule
\textbf{Apache Bench} (\texttt{ab}) & Un comando y una URL. \textbf{La de hoy.}\\
\texttt{hey} & Un comando, como \texttt{ab}\\
k6 & Un guion en JavaScript con varios pasos\\
Locust & Un guion en Python que simula usuarios\\
JMeter & Con ventanas, para pruebas grandes\\
\bottomrule
\end{tabular}
\end{columns}
\end{frame}

% 17
\begin{frame}[fragile]{Apache Bench: el comando}
\begin{lstlisting}[style=terminal,basicstyle=\ttfamily\small\color{white}]
ab -n 200 -c 20 https://cafe-api-123456789.us-central1.run.app/menu
\end{lstlisting}
\vspace{1mm}
\begin{columns}[T,onlytextwidth]
\column{.48\textwidth}
\renewcommand{\arraystretch}{1.25}
\small
\begin{tabular}{@{}>{\ttfamily}l>{\raggedright\arraybackslash}p{4.4cm}@{}}
\toprule
\textsf{\textbf{Parte}} & \textbf{Qué es}\\
\midrule
ab & Apache Bench, la herramienta\\
-n 200 & Pedidos \textbf{en total}\\
-c 20 & Pedidos \textbf{a la vez}: la concurrencia\\
…/menu & La URL, \textbf{con la ruta}\\
\bottomrule
\end{tabular}
\vspace{2mm}

{\scriptsize Viene en el paquete \texttt{apache2-utils}. En Cloud Shell, si \texttt{ab} no aparece:\\\texttt{sudo apt-get install -y apache2-utils}}
\column{.48\textwidth}
\begin{center}
\begin{tikzpicture}[font=\scriptsize]
  \node[caja,fill=utpGrisClaro,minimum width=1.2cm,minimum height=2.4cm] (ab) at (0,0) {\textbf{ab}};
  \node[caja,draw=okVerde,fill=okVerde!12,minimum width=1.4cm,minimum height=2.4cm] (api) at (4.2,0) {\textbf{API}};
  \foreach \y in {.9,.45,0,-.45,-.9} \draw[flecha] (ab.east |- 0,\y) -- (api.west |- 0,\y);
  \node[above=1pt of ab.north east,anchor=south west,font=\scriptsize] {20 a la vez};
  \node[text width=5.2cm,align=center,text=utpGris] at (2.1,-1.8) {cuando uno termina, sale otro,\\hasta completar 200};
\end{tikzpicture}
\end{center}
\end{columns}
\end{frame}

% 18
\begin{frame}[fragile]{Apache Bench: cómo leer el resultado}
\begin{columns}[T,onlytextwidth]
\column{.50\textwidth}
\begin{lstlisting}[style=terminal]
Concurrency Level:      20
Time taken for tests:   1.079 seconds
|Complete requests:      200|
|Failed requests:        0|
...
|Requests per second:    185.36 [#/sec]|
|Time per request:       107.90 [ms]|
...
Percentage of the requests served
within a certain time (ms)
  |50%    101|
  |95%    139|
 |100%    312 (longest request)|
\end{lstlisting}
{\scriptsize Un ejemplo: tus números serán otros.}
\column{.46\textwidth}
\renewcommand{\arraystretch}{1.2}
\small
\begin{tabular}{@{}>{\raggedright\arraybackslash}p{2.3cm}>{\raggedright\arraybackslash}p{3.7cm}@{}}
\toprule
\textbf{Línea} & \textbf{Qué te dice}\\
\midrule
\texttt{Complete} & 200 terminaron\\
\texttt{Failed} & \textcolor{okVerde}{\textbf{Ninguno falló}}\\
\texttt{Requests per second} & \textbf{185 pedidos por segundo}: el número principal\\
\texttt{Time per request} & Cada usuario esperó unos 108 ms\\
\texttt{95\%} & 95 de cada 100 tardaron 139 ms o menos\\
\texttt{100\%} & El más lento: 312 ms\\
\bottomrule
\end{tabular}
\end{columns}
\end{frame}

% 19
\begin{frame}{Cierre: preguntas para comprobar}
\begin{block}{Responde con tus palabras}
\small
1. Para App Engine subes \texttt{main.py}, \texttt{requirements.txt} y \texttt{app.yaml}. Para Cloud Run, ¿qué archivo cambia y qué tienes que construir antes?\\\medskip
2. En App Engine, ¿qué archivo dice qué Python usar? ¿Qué comando sube tu código?\\\medskip
3. Corres \texttt{ab -n 500 -c 50} y dice \texttt{Failed requests: 0} y \texttt{Requests per second: 210}. ¿Qué significan \texttt{-n}, \texttt{-c} y esos dos números?
\end{block}
\end{frame}

% 20
\begin{frame}{Referencias}
\scriptsize
\begin{itemize}\setlength\itemsep{4pt}
  \item App Engine: \texttt{cloud.google.com/appengine/docs/standard/python3}.
  \item Cloud Run: \texttt{cloud.google.com/run/docs}. Cloud Build y Artifact Registry: \texttt{cloud.google.com/build} y \texttt{cloud.google.com/artifact-registry}.
  \item Apache Bench: \texttt{httpd.apache.org/docs/2.4/programs/ab.html}.
  \item Otras herramientas de carga: \texttt{github.com/rakyll/hey}, \texttt{k6.io}, \texttt{locust.io} y \texttt{jmeter.apache.org}.
  \item Material previo del curso: Sesión 28 (la API de Café Nube y su Dockerfile) y Sesión 30 (Cloud Build, Artifact Registry y GKE).
  \item Las respuestas de las preguntas están en \texttt{../respuestas.md}.
\end{itemize}
\end{frame}

\end{document}
